\chapter{تطبيقات المؤسسة}
\label{ch:org}

تتعرف على التطبيقات الجاهزة التي تنظّم عمل المؤسسة نفسها: الموظفون والمناصب، والخطط والمشاريع، والكتب الرسمية، والإجراءات، وقواعد العمل، وكيف يعرف كل تطبيق ما سجّله الآخر.

\section{الموارد البشرية: الأساس الذي يبنى عليه كل شيء}
\label{sec:org-hr-base}
\begin{rcslide}{الموارد البشرية: الأساس}
\sphase{الموارد البشرية}
\slead{يُسجَّل الموظف «عضواً»، ويُربط بمنصبه عبر «تعيين». ومن هذا الربط تعرف التطبيقات الأخرى من أنت وأين تعمل.}
\schain{العضو ← التعيين ← المنصب ← الوحدة التنظيمية والموقع}
\scard{people}{الأعضاء}{ملف الموظف وبياناته.}
\scard{building}{المناصب}{الدور الذي يشغله الموظف في الهيكل.}
\scard{link}{التعيينات}{الرابط بين الموظف ومنصبه.}
\stake{حين يكون الهيكل صحيحاً، تعرف التطبيقات من يشغل كل منصب، فتُوجَّه الكتب إلى المنصب الصحيح دون سؤال.}
\slink{main}{الأعضاء}{ar/hr-members.html}
\slink{main}{المناصب}{ar/hr-positions.html}
\slink{main}{التعيينات}{ar/hr-assignments.html}
\end{rcslide}

\section{الموارد البشرية: يومك أنت}
\label{sec:org-hr-day}
\begin{rcslide}{الموارد البشرية: يومك أنت}
\sphase{الموارد البشرية}
\slead{ما يخصك من الموارد البشرية تراه بنفسك، دون أن تسأل أحداً.}
\scard{flow}{الورديات}{وقت عملك المحدد، وهو أساس احتساب حضورك.}
\scard{chart}{حضوري}{شهرك كما احتسبه النظام: بصماتك وإجمالياتك وأثرها في الراتب.}
\scard{paper}{الإجازات}{طلب بنوع وتاريخ، وحالته: قيد الانتظار أو معتمدة أو مرفوضة.}
\scard{sheet}{قسائم الرواتب}{أيام العمل والاستحقاقات والاستقطاعات وصافي الراتب.}
\stake{ما كنت تسأل عنه الموارد البشرية صار أمامك، والرقم نفسه يراه الموظف والمدير.}
\slink{main}{حضوري}{ar/hr-my-attendance.html}
\slink{main}{الإجازات}{ar/hr-leaves.html}
\slink{main}{الورديات}{ar/hr-shifts.html}
\end{rcslide}

\section{الاستراتيجية والتنفيذ}
\label{sec:org-projects}
\begin{rcslide}{الاستراتيجية والتنفيذ}
\sphase{الخطط والمشاريع}
\slead{للأعمال الكبيرة نتيجة مطلوبة ومسؤول وميزانية، وتُراجَع مخرجاتها قبل اعتبارها مكتملة.}
\schain{الخطة ← المشروع ← الميزانية ← المخرجات}
\scard{flow}{الخطة}{النتيجة المطلوبة والجهة المسؤولة والفترة.}
\scard{projects}{المشروع}{عمل محدد النطاق، له مدير وميثاق.}
\scard{chart}{الميزانية}{المخطط مقابل ما صُرف فعلاً.}
\scard{tasks}{المخرجات}{لكل مخرج موعد ومعيار قبول.}
\stake{تعرف في كل لحظة أين وصل المشروع، وكم صُرف منه، وما الذي سُلّم فعلاً.}
\slink{main}{المشاريع والتخطيط}{ar/c5.html}
\slink{main}{المشاريع}{ar/projects-projects.html}
\slink{main}{الميزانيات}{ar/projects-budgets.html}
\end{rcslide}

\section{المراسلات: صفتك وصندوقك}
\label{sec:org-corr-tray}
\begin{rcslide}{المراسلات: صفتك وصندوقك}
\sphase{المراسلات}
\slead{تعمل في المراسلات باسم منصبك، ويُسمّى «الصفة». ويجمع «صندوق المراسلات» عملك كله في مجلدات.}
\scard{people}{الصفة}{المنصب الذي تعمل باسمه. وتغييرها لا يمنحك صلاحية جديدة.}
\scard{mail}{صندوق المراسلات}{ما وصلك، وما ينتظر قرارك، وما أرسلته.}
\scard{paper}{المسودات}{تحفظ الكتاب، ثم تعود إليه قبل إرساله للمراجعة.}
\stake{تعمل باسم المنصب، فيعرف الجميع بأي صفة اتُّخذ كل قرار. وإذا فُوّضت عن منصب آخر، تميّز العمل نيابةً عنه من عمل منصبك.}
\slink{main}{الصفة والصلاحيات}{ar/correspondence-capacity.html}
\slink{main}{صندوق المراسلات ودورة العمل}{ar/correspondence-tray.html}
\end{rcslide}

\section{المراسلات: دورة الكتاب}
\label{sec:org-corr-cycle}
\begin{rcslide}{المراسلات: دورة الكتاب}
\sphase{المراسلات}
\slead{يمر الكتاب بمراحل واضحة، وكل قرار فيها يحدد صاحب الخطوة التالية.}
\schain{المسودة ← المراجعة والاعتماد ← الإصدار ← الإرسال ← التسليم}
\scard{tasks}{الإصدار}{يعطي الكتاب رقمه.}
\scard{mail}{الإرسال}{يسجّل خروج الكتاب.}
\scard{paper}{التسليم}{يوثّق وصوله باليد.}
\stake{لا يضيع كتاب بين المكاتب: تعرف في أي مرحلة هو، ومن ينتظر قراره الآن.}
\slink{correspondence}{تطبيق عملي: المراسلات خطوة بخطوة}{ar/index.html}
\slink{main}{الصادر: الترقيم والإرسال والتسليم}{ar/correspondence-outgoing-dispatch.html}
\end{rcslide}

\section{دليل الإجراءات}
\label{sec:org-procedures}
\begin{rcslide}{دليل الإجراءات}
\sphase{الإجراءات}
\slead{دليل إجراءات المؤسسة بطاقاتٍ حيّة لها مالك وتُقاس، بدل ملف PDF لا يقرؤه أحد.}
\scompare{دليل ورقي}{ملف طويل يُطبع مرة \\ لا يُعرف من يحدّثه \\ تقرؤه ثم تبدأ العمل في مكان آخر}{دليل الإجراءات في كيان}{بطاقة لكل إجراء، لها مالك \\ «اقتراح تعديل» على أي بطاقة \\ «بدء الإجراء» يفتح المعاملة مباشرة}
\scard{paper}{استيراد دليل PDF}{تبدأ من الدليل الذي تملكه المؤسسة اليوم.}
\scard{mail}{من كتاب وارد}{يبدأ الإجراء من كتاب وارد مسجّل.}
\stake{الإجراء لا يبقى نصاً تقرؤه؛ تبدؤه من بطاقته، فيصبح معاملة تتابعها حتى نهايتها.}
\end{rcslide}

\section{نتائج الذكاء الاصطناعي}
\label{sec:org-findings}
\begin{rcslide}{نتائج الذكاء الاصطناعي}
\sphase{الرقابة}
\slead{أمور ترصدها سياسات المؤسسة في السجلات والطلبات والمستندات، وتحتاج إلى قرار من شخص.}
\schain{سياسة ترصد ← ملاحظة ← معالجتها أو استبعادها بسبب ← أو تحويلها إلى مهمة}
\scard{ai}{سياسات يقترحها الذكاء الاصطناعي}{ويختبرها على بياناتك قبل تطبيقها.}
\scard{people}{القرار لشخص}{لا يعمل شيء حتى يطبّقه شخص.}
\stake{لا تنتظر أن يكتشف أحد الخطأ صدفة؛ تصل الملاحظة إلى من يقرر فيها، ويبقى أثر قراره محفوظاً.}
\end{rcslide}

\section{تطبيقات تظهر حين تحتاجها}
\label{sec:org-more}
\begin{rcslide}{تطبيقات تظهر حين تحتاجها}
\sphase{تطبيقات أخرى}
\slead{بعض التطبيقات لا تراها إلا إذا كان لك دور فيها، أو فعّلتها مؤسستك.}
\scard{form}{خطط التكليف}{من يعبّئ أي نموذج، وعن أي سجلات، وإلى أين وصلت كل جولة.}
\scard{tasks}{تقييماتي}{السجلات التي أنت في لجنة تقييمها، ومواعيدها.}
\scard{chat}{الإعلانات}{إعلانات المؤسسة إلى جمهور تحدده.}
\stake{إذا ذكر زميلك تطبيقاً لا تجده، فقد لا يكون لك دور فيه بعد.}
\end{rcslide}

\section{الإعدادات: قواعد المؤسسة}
\label{sec:org-settings}
\begin{rcslide}{الإعدادات: قواعد المؤسسة}
\sphase{الإعدادات}
\slead{يضبط مسؤول المؤسسة القواعد مرة واحدة، فتعمل بها التطبيقات كلها. وأنت ترى أثرها في عملك.}
\scard{building}{إعدادات المؤسسة}{الهوية والقيم الافتراضية.}
\scard{lock}{الأدوار}{ما يراه كل مستخدم وما يفعله.}
\scard{tasks}{النشر والاعتماد}{التغييرات التي تحتاج إلى اعتماد قبل تطبيقها.}
\scard{paper}{حفظ السجلات}{مدة بقاء السجلات ومصيرها.}
\stake{إذا لم تجد صفحة أو زراً تتوقعه، فالسبب غالباً في دورك. اسأل مسؤول المؤسسة.}
\slink{main}{الأدوار}{ar/settings-roles.html}
\slink{main}{الإعدادات}{ar/c8.html}
\end{rcslide}

\section{تطبيقات يعرف بعضها بعضاً}
\label{sec:org-linked}
\begin{rcslide}{تطبيقات يعرف بعضها بعضاً}
\sphase{الخلاصة}
\slead{ما تسجّله في تطبيق لا تعيد كتابته في غيره.}
\schain{ملفك في الموارد البشرية ← تعيينك في منصب ← صفتك في المراسلات}
\smetric{بيانات الموظف}{في كل برنامج نسخة}{تُسجَّل مرة واحدة}
\smetric{من المسؤول؟}{بحسب الذاكرة}{مكتوب في السجل}
\smetric{أين وصل العمل؟}{بالسؤال}{بالحالة}
\stake{هذا الترابط هو ما يجعل أي تقرير ممكناً من مكان واحد، كما رأيت في الوحدة الأولى.}
\end{rcslide}
